% !TEX program = lualatex
\documentclass[11pt,a4paper]{article}

\usepackage{luatexja}
\usepackage[haranoaji]{luatexja-preset}
\usepackage{amsmath,amssymb,mathtools}
\usepackage{booktabs}
\usepackage{geometry}
\geometry{top=25mm,bottom=25mm,left=25mm,right=25mm}
\usepackage{hyperref}
\usepackage{xcolor}

\title{Timefold Vehicle Routing (MDVRPTW) の解説と定式化}
\author{Timefold Quickstarts Use-Case Analysis}
\date{\today}

\begin{document}

\maketitle

\section{問題の概要}

Timefold QuickstartsのVehicle Routing問題は、複数拠点（デポ）、時間枠（Time Windows）および車両容量（Capacity）を考慮した複数デポ配送計画問題（\textbf{MDVRPTW: Multi-Depot Vehicle Routing Problem with Time Windows}）です。

指定された複数のデポ（拠点・車庫）を出発した複数台の車両が、それぞれの拠点を起点として訪問先（顧客）の需要を満たしつつ効率的に巡回し、再びそれぞれの拠点へ戻るルートを計画します。

Timefold内部では \texttt{HardMediumSoftScore} を用いた3階層のスコア評価が行われており、制約の分類は以下の通りです。

\begin{table}[htbp]
  \centering
  \caption{Timefold Vehicle Routing の制約一覧}
  \begin{tabular}{lll}
    \toprule
    制約名 & スコアレベル & 説明 \\
    \midrule
    \texttt{vehicleCapacity} & Hard & 各車両の訪問先の需要合計が積載容量を超えてはならない \\
    \texttt{serviceFinishedAfterMaxEndTime} & Hard & 訪問先でのサービス完了時刻が最大終了時刻を超えてはならない \\
    \texttt{maximizeVisitsAssigned} & Medium & 可能な限り多くの訪問先を車両に割り当てる \\
    \texttt{minimizeTravelTime} & Soft & 全車両の総移動時間（または総距離）を最小化する \\
    \bottomrule
  \end{tabular}
\end{table}

\section{記号の定義}

\subsection{集合およびインデックス}
\begin{itemize}
    \item $V = \{1, 2, \dots, n\}$ : 訪問先（顧客）の集合
    \item $D = \{d_1, d_2, \dots, d_m\}$ : デポ（拠点・車庫）の集合
    \item $N = V \cup D$ : 全ノードの集合
    \item $K = \{1, 2, \dots, m\}$ : 利用可能な車両の集合
\end{itemize}

\subsection{パラメータ（定数）}
\begin{itemize}
    \item $h(k) \in D$ : 車両 $k \in K$ に割り当てられた拠点デポ (\texttt{homeLocation})
    \item $d_i$ : 訪問先 $i \in V$ の需要量（$i \in D$ のとき $d_i = 0$）
    \item $C_k$ : 車両 $k \in K$ の積載容量
    \item $t_{ijk}$ : 車両 $k$ がノード $i$ からノード $j$ へ移動する時間（$i, j \in N$）
    \item $s_i$ : 訪問先 $i \in V$ でのサービス時間（$i \in D$ のとき $s_i = 0$）
    \item $[a_i, b_i]$ : 訪問先 $i \in V$ の時間枠（$a_i$: 開始可能時刻 \texttt{minStartTime}、$b_i$: 最大終了時刻 \texttt{maxEndTime}）
    \item $e_k$ : 車両 $k \in K$ のデポ出発時刻 (\texttt{departureTime})
    \item $M$ : 十分に大きな正の実数（Big-$M$）
    \item $W_{\text{medium}}, W_{\text{soft}}$ : 目的関数の重みパラメータ（$W_{\text{medium}} \gg W_{\text{soft}} > 0$）
\end{itemize}

\subsection{決定変数}
\begin{itemize}
    \item $x_{ijk} \in \{0, 1\}$ : 車両 $k$ がノード $i$ からノード $j$ へ直接移動するとき $1$、それ以外は $0$
    \item $S_i \ge 0$ : 訪問先 $i \in V$ でのサービス開始時刻
    \item $u_i \in \{0, 1\}$ : 訪問先 $i \in V$ がどの車両にも割り当てられない（未訪問）とき $1$、割当済のとき $0$
\end{itemize}

\section{数理モデルの定式化}

\subsection{目的関数}
未訪問ペナルティ（Mediumレベル）および総移動時間（Softレベル）をペナルティ項として最小化します。
\begin{equation}
\min \quad W_{\text{medium}} \sum_{i \in V} u_i \;+\; W_{\text{soft}} \sum_{k \in K} \sum_{i \in N} \sum_{\substack{j \in N \\ j \neq i}} t_{ijk} \, x_{ijk}
\end{equation}

\subsection{制約条件}

\subsubsection{1. 訪問割当およびフロー保存則}
各訪問先 $i \in V$ は高々1回訪問される：
\begin{equation}
\sum_{k \in K} \sum_{\substack{j \in N \\ j \neq i}} x_{ijk} = 1 - u_i \quad \forall i \in V
\end{equation}

各車両 $k$ および訪問先 $i \in V$ における流出入の一致：
\begin{equation}
\sum_{\substack{j \in N \\ j \neq i}} x_{ijk} - \sum_{\substack{j \in N \\ j \neq i}} x_{jik} = 0 \quad \forall i \in V, \, \forall k \in K
\end{equation}

各車両 $k$ は割り当てられた拠点 $h(k)$ を出発し、再び同じ拠点に戻る：
\begin{equation}
\sum_{j \in V} x_{h(k) j k} \le 1 \quad \forall k \in K
\end{equation}
\begin{equation}
\sum_{i \in V} x_{i h(k) k} = \sum_{j \in V} x_{h(k) j k} \quad \forall k \in K
\end{equation}

車両 $k$ は割り当てられていない他車両の拠点 $d \in D \setminus \{h(k)\}$ からの出入りを行わない：
\begin{equation}
\sum_{\substack{j \in N \\ j \neq d}} x_{d j k} = 0 \quad \forall k \in K, \, \forall d \in D \setminus \{h(k)\}
\end{equation}
\begin{equation}
\sum_{\substack{i \in N \\ i \neq d}} x_{i d k} = 0 \quad \forall k \in K, \, \forall d \in D \setminus \{h(k)\}
\end{equation}

\subsubsection{2. 車両容量制約 (Hard)}
車両 $k$ に割り当てられた全訪問先の需要合計は容量 $C_k$ を超えない：
\begin{equation}
\sum_{i \in V} d_i \left( \sum_{\substack{j \in N \\ j \neq i}} x_{ijk} \right) \le C_k \quad \forall k \in K
\end{equation}

\subsubsection{3. 時間枠・完了時刻制約およびサブツアー排除 (Hard)}
サービス開始時刻の推移（部分巡回路の排除を兼ねる）：
\begin{equation}
\tau_{ik} + t_{ijk} - M (1 - x_{ijk}) \le S_j \quad \forall i \in N, \, \forall j \in V, \, i \neq j, \, \forall k \in K
\end{equation}
ただし、出発元の時刻 $\tau_{ik}$ は以下のように定義される：
\begin{equation}
\tau_{ik} = 
\begin{cases}
S_i + s_i & (i \in V \text{ の場合}) \\
e_k & (i \in D \text{ の場合})
\end{cases}
\end{equation}

最小開始可能時刻以降にサービスを開始する：
\begin{equation}
S_i \ge a_i \quad \forall i \in V
\end{equation}

サービス完了時刻が最大終了時刻を超えてはならない：
\begin{equation}
S_i + s_i \le b_i \quad \forall i \in V
\end{equation}

\subsubsection{4. 変数の定義域}
\begin{align}
x_{ijk} &\in \{0, 1\} \quad \forall i, j \in N, \, \forall k \in K \\
u_i &\in \{0, 1\} \quad \forall i \in V \\
S_i &\ge 0 \quad \forall i \in V
\end{align}

\end{document}

